\documentclass[11pt,a4paper]{article}
\usepackage[margin=1in]{geometry}
\usepackage{times}
\usepackage{amsmath}
\usepackage{enumitem}
\usepackage{booktabs}
\usepackage{array}
\usepackage[colorlinks,allcolors=blue]{hyperref}
\setlist{leftmargin=1.4em,topsep=2pt,itemsep=2pt}
\setlength{\parindent}{0pt}
\setlength{\parskip}{4pt}

\title{\vspace{-1.5em}Initial Project Plan\\[2pt]
\large LSTM vs.\ Transformer Encoder for Text Classification (AG News)}
\author{Jihun Jang (ETRI School, 02521122) \quad Deep Learning 2, 2026}
\date{}

\begin{document}
\maketitle
\vspace{-1.5em}

\section*{1.\ Team}
\begin{itemize}
\item Team: Jihun Jang (ETRI School, 02521122).
\item Team leader: Jihun Jang (solo team).
\end{itemize}

\section*{2.\ Dataset Split}
\begin{itemize}
\item HuggingFace \texttt{ag\_news}, single source, 4 classes (labels 0--3, no conversion needed).
\item \texttt{train\_test\_split(test\_size=0.1, seed=42)} $\rightarrow$ train 108k / val 12k / test 7.6k. Classes balanced.
\item The test set is used for final evaluation only (no model selection, tuning, or early stopping).
\end{itemize}

\section*{3.\ Preprocessing}
\begin{itemize}
\item Tokenization: word level (lowercased, punctuation split).
\item Vocabulary: built from train only, max 20k, minimum frequency 2, with PAD and UNK tokens.
\item Maximum length 128 (truncate if longer, pad if shorter); padding is masked.
\item The same input is used for both models.
\end{itemize}

\section*{4.\ Model Configuration}
\begin{center}
\begin{tabular}{>{\raggedright\arraybackslash}p{2.4cm} >{\raggedright\arraybackslash}p{4.5cm} >{\raggedright\arraybackslash}p{6.5cm}}
\toprule
 & \textbf{LSTM} & \textbf{Transformer Encoder} \\
\midrule
Embedding & 128 & 128 (= $d_{\text{model}}$) \\
Core & BiLSTM, 2 layers, hidden 128 & 2 layers, 4 attention heads, FFN 256, sinusoidal positional encoding \\
Pooling & masked mean & masked mean \\
Output / Loss & Linear $\rightarrow$ 4, Cross-Entropy & Linear $\rightarrow$ 4, Cross-Entropy \\
\bottomrule
\end{tabular}
\end{center}

\begin{itemize}
\item Common: Adam, learning rate 1e-3, batch 64, up to 8 epochs, dropout 0.1 (common start, tuned by validation), seed 42, model selection by validation only.
\item Parameter counts: LSTM $\approx$ 3.22M / Transformer $\approx$ 2.83M (comparable scale). No pretrained models.
\end{itemize}

\section*{5.\ Ablation Plan}
\begin{itemize}
\item Topic: data efficiency --- train 5 / 25 / 50 / 100\% ($\approx$ 5k / 27k / 54k / 108k).
\item Hypothesis: with little data the Transformer is disadvantaged (weaker inductive bias) while the LSTM is stable on small data; the gap shrinks as data grows.
\item Controls: the vocabulary is built once on the full training set and fixed/reused; subsets keep class proportions (stratified) with a fixed seed; val/test always use the full set.
\item Figure: $x$ = data size, $y$ = accuracy / macro-F1, both models on one graph; the 100\% point is reused as the main comparison.
\item Optional: Transformer internal analysis --- positional encoding on/off, attention-weight visualization, gradient saliency (input attribution).
\end{itemize}

\section*{6.\ Expected Analysis}
\begin{itemize}
\item Performance: both around 90\%+, with the Transformer slightly ahead or comparable.
\item Data efficiency: at 5\% LSTM $\geq$ Transformer, at 100\% comparable (hypothesis test).
\item Convergence: the Transformer is fast but learning-rate sensitive; the LSTM is stable (checked via loss curves).
\item Confusion: Business $\leftrightarrow$ Sci/Tech expected to be largest. Analyze $\geq$ 5 misclassifications (shared / LSTM-only / Transformer-only): ambiguous wording, short sentences, truncation.
\item Mechanism: accuracy maintained with positional encoding off would indicate AG News is vocabulary-centric (bag-of-words); attention and saliency should concentrate on class keywords.
\end{itemize}

\vspace{0.5em}\hrule\vspace{0.5em}
\textbf{Environment:} Mac MPS / Colab T4, seed 42, checkpoints to Drive. \quad
\textbf{Deliverables:} Jupyter notebook + report (4--6 pp.) + AI Usage Appendix ($\leq$ 1 p.).

\end{document}
